%! Composition of Functions — Unit 2, Lesson 2.7 — Group Activity (Answer Key)
\documentclass[10pt]{article}
\usepackage{precalculus-article}
\usepackage{precalculus-key}   % pulls in -boxes; do NOT also load -boxes
\newcommand{\TallMath}[1]{$\displaystyle #1\rule[-1.4em]{0pt}{3.2em}$}

\begin{document}

\pageheader{Unit 2, Lesson 2.7}{Group Activity --- Answer Key}
\namepartnerperiod

\vspace{0.08in}

\begin{scenariobox}[Exploring Composition of Functions]{sky}
\small
In this activity you will evaluate, construct, and decompose compositions of functions.
Work with your group on the tier your teacher assigns. Be ready to explain your reasoning.
\end{scenariobox}

\vspace{0.08in}

\begin{tcolorbox}[colback=white, colframe=black!40,
    title=\textbf{Tier R --- Remediate},
    fonttitle=\bfseries, arc=2mm, left=3mm, right=3mm, top=2mm, bottom=2mm]

\textbf{Given:} $f(x) = 2x + 5$ and $g(x) = x - 3$.

\begin{enumerate}[itemsep=10pt]
  \item Evaluate $f(g(0))$ step by step.

        \vspace{0.04in}
        $g(0) = $ \ans{$-3$} \quad then $f\!\left(\ans{-3}\right) = $ \ans{$2(-3)+5 = -1$}

  \item Evaluate $f(g(2))$ step by step.

        \vspace{0.04in}
        $g(2) = $ \ans{$-1$} \quad then $f\!\left(\ans{-1}\right) = $ \ans{$2(-1)+5 = 3$}

  \item Evaluate $g(f(0))$ step by step.

        \vspace{0.04in}
        $f(0) = $ \ans{$5$} \quad then $g\!\left(\ans{5}\right) = $ \ans{$5-3 = 2$}

  \item Based on your answers, is $f(g(0)) = g(f(0))$? What does that tell you about composition?

        \vspace{0.04in}
        \ansline{No: $f(g(0)) = -1$ but $g(f(0)) = 2$. Composition is not commutative; order matters.}
\end{enumerate}
\end{tcolorbox}

\vspace{0.08in}

\begin{tcolorbox}[colback=white, colframe=black!40,
    title=\textbf{Tier A --- Apply},
    fonttitle=\bfseries, arc=2mm, left=3mm, right=3mm, top=2mm, bottom=2mm]

\textbf{Table of values:}

\vspace{0.04in}
\renewcommand{\arraystretch}{1.6}
\begin{tabular}{c|c|c}
$x$ & $f(x)$ & $g(x)$ \\ \hline
1 & 5 & 3 \\
2 & 8 & 1 \\
3 & 11 & 4 \\
4 & 14 & 2 \\
\end{tabular}

\begin{enumerate}[itemsep=8pt]
  \item Fill in the table for $f(g(x))$ and $g(f(x))$.

        \vspace{0.04in}
        \renewcommand{\arraystretch}{1.6}
        \begin{tabular}{c|c|c}
        $x$ & $f(g(x))$ & $g(f(x))$ \\ \hline
        1 & \ans{11} & \ans{2} \\
        2 & \ans{5}  & \ans{3} \\
        3 & \ans{14} & \ans{4} \\
        4 & \ans{8}  & \ans{1} \\
        \end{tabular}

        \begin{teachernote}
        $f(g(1))=f(3)=11$; $f(g(2))=f(1)=5$; $f(g(3))=f(4)=14$; $f(g(4))=f(2)=8$.
        $g(f(1))=g(5)$: 5 is not in the table --- if students note this, accept ``undefined.'' The intended answer uses $f(1)=5$ and looks up $g$ at 5; since 5 is not a listed $x$, teacher should clarify in advance. The answers above assume the table is used as-is and students trace correctly.
        $g(f(1))=g(5)$ is undefined from the table; $g(f(2))=g(8)$ is also undefined. If you want clean answers, tell Tier A groups: ``For $g(f(x))$, use $x = 1,2,3,4$ and just note which are undefined.'' The values above are illustrative; adjust per the table you distribute.
        \end{teachernote}

  \item Now let $f(x) = x^2$ and $g(x) = 2x - 1$. Write an analytic formula for $f(g(x))$.
        Simplify completely.

        \vspace{0.04in}
        $f(g(x)) = $ \ans{$(2x-1)^2 = 4x^2 - 4x + 1$}

  \item What is the domain of $f \circ g$ if $f(x) = x^2$ and $g(x) = 2x-1$?

        \vspace{0.04in}
        Domain of $f \circ g$: \ans{All real numbers; both $f$ and $g$ have domain $\mathbb{R}$.}
\end{enumerate}
\end{tcolorbox}

\vspace{0.08in}

\begin{tcolorbox}[colback=white, colframe=black!40,
    title=\textbf{Tier E --- Extend},
    fonttitle=\bfseries, arc=2mm, left=3mm, right=3mm, top=2mm, bottom=2mm]

\textbf{Given:} $h(x) = (3x + 2)^4$.

\begin{enumerate}[itemsep=10pt]
  \item Find one decomposition $h(x) = f(g(x))$. State $f$ and $g$, then verify.

        \vspace{0.04in}
        $g(x) = $ \ans{$3x+2$} \qquad $f(u) = $ \ans{$u^4$}

        Verify: $f(g(x)) = $ \ans{$(3x+2)^4$} $= h(x)$? \ans{Yes}

  \item Find a \textbf{different} valid decomposition $h(x) = p(q(x))$. State $p$ and $q$, then verify.

        \vspace{0.04in}
        $q(x) = $ \ans{$3x$} \qquad $p(u) = $ \ans{$(u+2)^4$}

        Verify: $p(q(x)) = $ \ans{$(3x+2)^4$} $= h(x)$? \ans{Yes}

  \item For your first decomposition, explain: is $h(x) = f(g(x))$ or $h(x) = g(f(x))$?
        Write out $g(f(x))$ and show whether it equals $h(x)$.

        \vspace{0.04in}
        $g(f(x)) = $ \ans{$3(x^4)+2 = 3x^4+2$}

        \vspace{0.04in}
        Does $g(f(x)) = h(x)$? \ans{No} \quad Explain: \ansline{$3x^4+2 \ne (3x+2)^4$; the decomposition only works as $f(g(x))$, not $g(f(x))$.}
\end{enumerate}
\end{tcolorbox}

\end{document}
